% Grenzen + mögliche Weiterentwicklungen
\section{Grenzen des entwickelten Prototyps und mögliche Weiterentwicklungen}

Da die Dash\-board"=Lösung bislang ein Projekt prototypischen Charakters ist, kann dieses noch stark ausgeweitet werden. Neben dem Ausbau auf den restlichen Testpool kann das Nutzerfeedback zur Optimierung umgesetzt werden. Dieses liefert hilfreiche Hinweise bezüglich des bestehenden Arbeitsablaufs sowie der Formatierung, Größe und Kapselung zusammengehöriger Informationen, um den Übergang auf das neue System zu erleichtern.

Bisher bietet das Dash\-board rein inhaltlich betrachtet keinen Mehrwert. Da die Daten unabhängig vom bestehenden System ausgewertet und aufgearbeitet werden, sind vorgeschlagene Funktionen wie die Darstellung der Fehlerwerte der letzten zwei Wochen in einem Verlaufsdiagramm möglich. Da das Dash\-board jeden Tag überwacht wird, ist ein direkter Vergleichswert zum Vortag denkbar. Dieser zeigt zusätzlich zur Fehlerzahl, ob und wie viele Fehler mehr oder weniger als am Vortag vorliegen.

Da das System keine Veränderung am Altsystem vornimmt, kann diese Lösung ausschließlich als erster Schritt des Modernisierungsprozesses angesehen werden. Es hilft, Altlasten zu erkennen und fachfremden Personen den Einstieg in das System zu erleichtern. Ein nächster Schritt wäre, diese Erkenntnisse zur Umstrukturierung und Optimierung oder zur Umstrukturierung in Kombination mit einem Rehosting auf zum Beispiel ein standardisiertes, modernes Framework zu nutzen. Somit trägt diese Lösung aktiv zum nachhaltigen Bestand der Software bei, dessen Daten auf dem Dash\-board angezeigt werden.
